\section*{Capítulo \ref{ch:b1:diffeq} --- Equações diferenciais lineares}

\begin{solution}{exo:b1:diffeq:1}
Homogêneo: $y_h = \lambda\,\eu^{-2x}$. Particular: tente $ y_p =
c\,\eu^{3x}$ ($3$ não é uma raiz de $ r + 2$): $3c + 2c = 1$, $ c =
\frac15$. Solução geral: $ y = \frac{\eu^{3x}}{5} +
\lambda\,\eu^{-2x}$. Com $ y(0) = 1$: $\frac15 + \lambda = 1$,
$\lambda = \frac45$, então $ y = \frac{\eu^{3x} + 4\,\eu^{-2x}}{5}$.
\end{solution}

\begin{solution}{exo:b1:diffeq:2}
Sobre $\intoo{0}{+\infty}$, divida por $ x $: $ y' - \frac{2}{x}\,y = x^2$.
Aqui $ A(x) = -2\ln x $, $\eu^{-A(x)} = x^2$: soluções homogêneas
$\lambda x^2$. Variação de constantes: $\mu'(x) = x^2 \cdot x^{-2} =
1$, então $\mu = x + \lambda $ e
\[
y(x) = x^3 + \lambda x^2, \qquad \lambda \in \R .
\]
\emph{Verificar:} $ x(3x^2 + 2\lambda x) - 2(x^3 + \lambda x^2) = x^3$.
\end{solution}

\begin{solution}{exo:b1:diffeq:3}
$ y'' - 3y' + 2y = 0$: raízes $1$ e $2$; $\;y = \lambda\,\eu^{x} +
\mu\,\eu^{2x}$.

$ y'' + 4y' + 4y = 0$: raiz dupla $-2$; $\;y = (\lambda + \mu
x)\,\eu^{-2x}$.

$ y'' - 2y' + 5y = 0$: raízes $1 \pm 2\iu $; $\;y = \eu^{x}(\lambda\cos
2x + \mu\sin 2x)$.
\end{solution}

\begin{solution}{exo:b1:diffeq:4}
Homogêneo: raízes $\pm 1$, $ y_h = \lambda\,\eu^x + \mu\,\eu^{-x}$.
Particular com polinômio do lado direito ($\gamma = 0$ não é uma raiz):
$ y_p = ax^2 + bx + c $; substituindo, $2a - (ax^2 + bx + c) = x^2$
dá $ a = -1$, $ b = 0$, $ c = 2a = -2$: $ y_p = -x^2 - 2$. Geral
solução $ y = -x^2 - 2 + \lambda\eu^x + \mu\eu^{-x}$.

Cauchy: $ y(0) = -2 + \lambda + \mu = 0$ e $ y'(0) = \lambda - \mu =
1$: $\lambda = \frac32$, $\mu = \frac12$. Então
$ y = -x^2 - 2 + \frac{3\eu^x + \eu^{-x}}{2}$.
\end{solution}

\begin{solution}{exo:b1:diffeq:5}
$ a(x) = \tan x $, $ A(x) = -\ln(\cos x)$ (válido: $\cos > 0$ no
intervalo), $\eu^{-A} = \cos x $: soluções homogêneas $\lambda\cos x $.
Variação de constantes: $\mu'(x) = \sin 2x \cdot \frac{1}{\cos x} =
2\sin x $, então $\mu = -2\cos x + \lambda $ e
\[
y(x) = -2\cos^2 x + \lambda \cos x .
\]
\emph{Verificar:} $ y' = 4\cos x \sin x - \lambda\sin x $ e $ y\tan x =
-2\cos x\sin x + \lambda \sin x $; a soma deles é $2\cos x\sin x =
\sin 2x $, conforme necessário.
\end{solution}

\begin{solution}{exo:b1:diffeq:6}
$\chi(r) = r^2 - 4r + 3 = (r-1)(r-3)$: $\gamma = 1$ é uma raiz simples
($ m = 1$). Tentar $ y_p = x(ax + b)\,\eu^x $. Com $ u = ax^2 + bx $,
\[
y_p'' - 4y_p' + 3y_p = \bigl(u'' + (2 - 4)u' + \chi(1) u\bigr)\eu^x
= \bigl(2a - 2(2ax + b)\bigr)\eu^x .
\]
Identifique-se com $(2x + 1)\eu^x $: $-4a = 2$ e $2a - 2b = 1$, então $ a =
-\frac12$, $ b = -1$. Solução geral:
\[
y = -\Bigl(\frac{x^2}{2} + x\Bigr)\eu^{x} + \lambda\,\eu^{x} +
\mu\,\eu^{3x}, \qquad (\lambda, \mu) \in \R^2 .
\]
\end{solution}

\begin{solution}{exo:b1:diffeq:7}
Homogêneo: $ y_h = \lambda\cos 2x + \mu\sin 2x $.

Lado direito $ x $ ($\gamma = 0$ não é uma raiz): $ y_1 = ax + b $ com
$4(ax + b) = x $: $ y_1 = \frac x4$.

Lado direito $\sin 2x = \Im(\eu^{2\iu x})$, $2\iu $ raiz simples de
$ r^2 + 4$: tentar $ z = c\,x\,\eu^{2\iu x}$; então $ z'' + 4z =
4\iu c\,\eu^{2\iu x}$, igual a $\eu^{2\iu x}$ para $ c =
\frac{1}{4\iu} = -\frac{\iu}{4}$. Então $ z = -\frac{\iu x}{4}(\cos 2x +
\iu \sin 2x)$ e $ y_2 = \Im(z) = -\frac{x\cos 2x}{4}$.

Por superposição:
\[
y = \frac{x}{4} - \frac{x\cos 2x}{4} + \lambda\cos 2x + \mu\sin 2x .
\]
\end{solution}

\begin{solution}{exo:b1:diffeq:8}
A equação $ T' + kT = 20k $ tem solução particular constante $20$ e
soluções homogêneas $\lambda\eu^{-kt}$: $ T(t) = 20 +
\lambda\,\eu^{-kt}$, e $ T(0) = 80$ dá $\lambda = 60$:
\[
T(t) = 20 + 60\,\eu^{-kt} .
\]
$ T(10) = 50$: $\eu^{-10k} = \frac12$, então $ k = \frac{\ln 2}{10}$. Então
$ T(t) = 25$ requer $\eu^{-kt} = \frac{5}{60} = \frac{1}{12}$, ou seja
\[
t = \frac{\ln 12}{k} = 10\,\frac{\ln 12}{\ln 2} \approx 35.8
\text{ minutes.}
\]
\end{solution}

\begin{solution}{exo:b1:diffeq:9}
\begin{enumerate}
  \item $\chi(r) = r^2 + 2\varepsilon r + 1$, $\Delta = 4(\varepsilon^2
        - 1)$.
        Para $\varepsilon \in \intco{0}{1}$: raízes $-\varepsilon \pm
        \iu\sqrt{1 - \varepsilon^2}$, então
        $ y = \eu^{-\varepsilon t}\bigl(\lambda\cos\omega t +
        \mu\sin\omega t\bigr)$ com $\omega = \sqrt{1 -
        \varepsilon^2}$.
        Para $\varepsilon = 1$: raiz dupla $-1$, $ y = (\lambda + \mu
        t)\,\eu^{-t}$.
        Para $\varepsilon > 1$: raízes reais $ r_\pm = -\varepsilon \pm
        \sqrt{\varepsilon^2 - 1}$, ambos $< 0$, e $ y =
        \lambda\eu^{r_+t} + \mu\eu^{r_-t}$.
  \item Para $\varepsilon \in \intoo{0}{1}$: $\abs y \leq
        \eu^{-\varepsilon t}(\abs\lambda + \abs\mu) \to 0$. Para
        $\varepsilon = 1$: $(\lambda + \mu t)\eu^{-t} \to 0$
        (exponencial vence polinômio,
        \cref{prop:b1:functions:powerrules}). Para $\varepsilon > 1$:
        ambas as exponenciais decaem desde $ r_\pm < 0$ (de fato
        $\sqrt{\varepsilon^2 - 1} < \varepsilon $). Para $\varepsilon =
        0$: $ y = \lambda\cos t + \mu\sin t $ tem amplitude constante
        $\sqrt{\lambda^2 + \mu^2} \neq 0$ a menos que $ y = 0$.
  \item Escrever $\lambda\cos\omega t + \mu\sin\omega t =
        R\cos(\omega t - \varphi)$ com $ R = \sqrt{\lambda^2 + \mu^2}
        > 0$. Os zeros de $ y $ são aqueles de $\cos(\omega t -
        \varphi)$ (o fator $\eu^{-\varepsilon t}$ nunca
        desaparece): $\omega t - \varphi \equiv \frac\pi2 \pmod \pi $,
        uma progressão aritmética com lacuna $\frac{\pi}{\omega} =
        \frac{\pi}{\sqrt{1 - \varepsilon^2}}$.
\end{enumerate}
\end{solution}

\begin{solution}{exo:b1:diffeq:10}
\omterm{def:b1:logic:sets}{Definir} $ x = y = 0$: $2f(0) = 2f(0)^2$, e $ f(0) \neq 0$ dá $ f(0) =
1$. Consertar $ x $ e diferencie a equação duas vezes em relação a $ y $:
\[
f''(x+y) + f''(x-y) = 2 f(x) f''(y) .
\]
Contexto $ y = 0$: $\;2f''(x) = 2 f(x) f''(0)$, aquilo é
\[
f''(x) = c\,f(x), \qquad c = f''(0).
\]
\emph{Caso $ c = \omega^2 > 0$:} $ f(x) = \lambda\cosh\omega x +
\mu\sinh\omega x $; $ f(0) = 1$ dá $\lambda = 1$. Conectando-se ao
equação funcional e usando as fórmulas de adição
(\cref{prop:b1:functions:hyprules}), a equação força $\mu = 0$
(compare os coeficientes de $\sinh\omega x \sinh\omega y $ ou
avaliar em $ x = y $): $ f = \cosh\omega x $, o que satisfaz
$\cosh(x+y) + \cosh(x-y) = 2\cosh x\cosh y $.

\emph{Caso $ c = -\omega^2 < 0$:} de forma similar $ f(x) = \cos\omega x $
($\omega \neq 0$), o que satisfaz a equação.

\emph{Caso $ c = 0$:} $ f $ afim com $ f(0) = 1$: $ f(x) = 1 + \mu x $;
as forças da equação $\mu = 0$, excluído ($ f $ não constante).

Conclusão: as soluções são $ f(x) = \cos\omega x $ e $ f(x) =
\cosh\omega x $, $\omega > 0$.
\end{solution}

\begin{solution}{exo:b1:diffeq:11}
\omterm{def:b1:logic:sets}{Definir} $ z(t) = y(\eu^t)$, para que $ y(x) = z(\ln x)$ para $ x > 0$. Então
\[
y'(x) = \frac{z'(\ln x)}x,
\qquad
y''(x) = \frac{z''(\ln x) - z'(\ln x)}{x^2} ,
\]
e substituindo na equação:
\[
x^2y'' - xy' + y = \bigl(z'' - z'\bigr) - z' + z
= z'' - 2z' + z = 0 .
\]
Polinômio característico $(r - 1)^2$: raiz dupla $1$, então $ z(t) =
(\lambda + \mu t)\,\eu^t $ e, de volta à variável $ x = \eu^t $:
\[
y(x) = (\lambda + \mu\ln x)\,x,
\qquad \lambda, \mu \in \R .
\]
\end{solution}

\begin{solution}{exo:b1:diffeq:12}
\begin{enumerate}
  \item Na forma normalizada $ y' - \frac2x\,y = 0$ em cada intervalo:
        $ A(x) = -2\ln\abs x $, então as soluções são $ y = a x^2$ sobre
        $\intoo0{+\infty}$ e $ y = b x^2$ sobre $\intoo{-\infty}0$,
        com constantes independentes
        (\cref{thm:b1:diffeq:homogeneous1}).
  \item Deixar $ y = ax^2$ para $ x \geq 0$ e $ bx^2$ para $ x < 0$. Ligado
        cada meia linha aberta $ y $ é diferenciável com $ xy' = 2y $.
        No $0$: os quocientes de diferença $\frac{y(h) - y(0)}h =
        ah $ ou $ bh $ tendem a $0$, então $ y'(0) = 0$ existe, e o
        equação em $ x = 0$ lê $0 \cdot y'(0) = 2y(0) = 0$:
        satisfeito. Então $ y $ resolve a equação de todos $\R $.
  \item As soluções em $\R $ são exatamente essas funções coladas: a
        \emph{dois}família de parâmetros para uma equação de primeira ordem.
        Não há contradição com \cref{thm:b1:diffeq:voc},
        cujas hipóteses falham aqui: escritas como $ y' + a(x)y = 0$,
        o coeficiente $ a(x) = -\frac2x $ não é contínuo em $0$
        --- na verdade não definido --- então $\R $ não é um intervalo em
        qual o teorema se aplica. A singularidade em $0$
        desconecta as duas meias-linhas, e o valor $ y(0) = 0$
        é forçado, não transportando nenhuma informação. Cada Cauchy
        dado em $ x_0 \neq 0$ determina a solução apenas no
        meia linha contendo $ x_0$.
\end{enumerate}
\end{solution}

\begin{solution}{pb:b1:diffeq:1}
\textbf{1.} $\chi(r) = r^2 + 2\lambda r + \omega_0^2$, $\Delta =
4(\lambda^2 - \omega_0^2) < 0$ para $0 < \lambda < \omega_0$: raízes
$-\lambda \pm \iu\omega_d $, então por
\cref{thm:b1:diffeq:homogeneous2}
\[
x(t) = \eu^{-\lambda t}\bigl(\lambda_1\cos\omega_d t +
\mu_1\sin\omega_d t\bigr), \qquad (\lambda_1, \mu_1) \in \R^2 .
\]
Para $\lambda = 0$: $ x = \lambda_1\cos\omega_0 t +
\mu_1\sin\omega_0 t $. A crítica ($\lambda = \omega_0$) e
superamortecido ($\lambda > \omega_0$) regimes são aqueles de
\cref{exo:b1:diffeq:9} (após o tempo de redimensionamento): $(\lambda_1 +
\mu_1 t)\eu^{-\lambda t}$, resp.\ combinações de $\eu^{r_\pm t}$
com $ r_\pm = -\lambda \pm \sqrt{\lambda^2 - \omega_0^2}$.

\textbf{2.} Subamortecido: $\abs x \leq \eu^{-\lambda
t}(\abs{\lambda_1} + \abs{\mu_1}) \to 0$. Crítico: $(\lambda_1 +
\mu_1 t)\eu^{-\lambda t} \to 0$ desde que as exponenciais bateram
polinômios (\cref{prop:b1:functions:powerrules}). Sobreamortecido:
$ r_- < r_+ = -\lambda + \sqrt{\lambda^2 - \omega_0^2} < 0$
porque $\sqrt{\lambda^2 - \omega_0^2} < \lambda $; ambos
decadência exponencial.

\textbf{3.} Ao longo de uma solução de $(H)$, usando $ x'' = -2\lambda x' -
\omega_0^2 x $:
\[
\mathcal E'(t) = x'x'' + \omega_0^2 x x'
= x'\bigl(-2\lambda x' - \omega_0^2 x\bigr) + \omega_0^2 xx'
= -2\lambda\,x'^2 \leq 0 .
\]
Se $ x(t_0) = x'(t_0) = 0$ então $\mathcal E(t_0) = 0$; $\mathcal E $
é não negativo e não crescente, então $\mathcal E \equiv 0$ sobre
$\intco{t_0}{+\infty}$, forçando $ x \equiv 0$ lá; para $ t \leq
t_0$, execute o mesmo argumento em $\tilde x(t) = x(2t_0 - t)$, que
resolve a equação com amortecimento $-\lambda $ mas ainda tem
$\tilde{\mathcal E}(t_0) = 0$ e $\tilde{\mathcal E}' =
+2\lambda\tilde x'^2 \geq 0$ com $\tilde{\mathcal E} \geq 0$;
não negativo, não decrescente e zero na extremidade direita de
$\intoc{-\infty}{t_0}$ significa zero por toda parte. Então $ x \equiv 0$ sobre
$\R $ --- uma prova energética de unicidade, válida para todos $\lambda
\geq 0$.

\textbf{4.} $ x(t + T_d) = R\,\eu^{-\lambda t}\eu^{-\lambda
T_d}\cos(\omega_d t + 2\pi - \varphi) = \eu^{-\lambda T_d}x(t)$.
O fator de redução por pseudoperíodo é $\eu^{-\delta}$ com
$\delta = \lambda T_d = \frac{2\pi\lambda}{\omega_d}$. Para
$\omega_0 = 1$, $\lambda = 0.1$: $\omega_d = \sqrt{0.99} =
0.99499$, então $\delta = \frac{0.62832}{0.99499} = 0.6315$: cada
balanço mantém $\eu^{-0.63} \approx 53\%$ da sua amplitude.

\textbf{5.} O fator de amplitude é $\eu^{-\lambda t}$, que
decai por $\eu $ sobre $ t = \frac1\lambda $. Esse intervalo contém
$\frac{1/\lambda}{T_d} = \frac{\omega_d}{2\pi\lambda}$
pseudo-períodos. Para $\lambda \ll \omega_0$, $\omega_d \approx
\omega_0$ e isso é $\approx \frac{\omega_0}{2\pi\lambda} =
\frac Q\pi $. Um $ Q = 300$ anéis de cordas de violão por cerca de cem
períodos; um $ Q = 1$ o amortecedor da porta não completa um.

\textbf{6.} Substituindo $ z\,\eu^{\iu\Omega t}$ para a mão esquerda
lado dá $ z\,(-\Omega^2 + 2\iu\lambda\Omega +
\omega_0^2)\,\eu^{\iu\Omega t}$, que é igual $ A\,\eu^{\iu\Omega
t}$ exatamente para $ z = \frac{A}{\omega_0^2 - \Omega^2 +
2\iu\lambda\Omega}$ (o denominador é diferente de zero: seu imaginário
parte é $2\lambda\Omega > 0$). Como os coeficientes são reais,
a verdadeira parte $ x_p = \Re\bigl(z\eu^{\iu\Omega t}\bigr)$ resolve o
equação com lado direito $\Re\bigl(A\eu^{\iu\Omega t}\bigr) =
A\cos\Omega t $.

\textbf{7.} Escrever $\omega_0^2 - \Omega^2 + 2\iu\lambda\Omega =
\sqrt{D}\,\eu^{\iu\varphi}$ com $ D = (\omega_0^2 - \Omega^2)^2 +
4\lambda^2\Omega^2$ e $\varphi \in \intoo0\pi $ (o imaginário
parte $2\lambda\Omega $ é positivo), então $\tan\varphi =
\frac{2\lambda\Omega}{\omega_0^2 - \Omega^2}$. Então $ z =
\frac{A}{\sqrt D}\eu^{-\iu\varphi}$ e
\[
x_p(t) = \Re\Bigl(\frac A{\sqrt D}\,\eu^{\iu(\Omega t -
\varphi)}\Bigr) = R\cos(\Omega t - \varphi),
\qquad R = \frac A{\sqrt D} .
\]

\textbf{8.} Como $\Omega \to 0^+$: $ D \to \omega_0^4$, então $ R \to
A/\omega_0^2$ e $\tan\varphi \to 0^+$ com $\varphi \in
\intoo0{\frac\pi2}$: $\varphi \to 0$. A massa segue a força
quase estaticamente, deslocado pela força/rigidez. Como $\Omega \to
+\infty $: $ D \sim \Omega^4$, então $ R \sim A/\Omega^2 \to 0$, e
$\varphi \to \pi $ (o número complexo $\omega_0^2 - \Omega^2 +
2\iu\lambda\Omega $ vai para o segundo quadrante com argumento $\to
\pi $): a massa mal se move e em oposição à fase ---
a inércia domina.

\textbf{9.} Por \cref{thm:b1:diffeq:structure2}~(1), cada solução
é $ x = x_p + x_h $ com $ x_h $ resolvendo $(H)$; pela pergunta 2, $ x_h(t)
\to 0$, então $ x(t) - x_p(t) \to 0$: todas as soluções convergem para o
mesmo estado estacionário. As condições iniciais apenas moldam o transitório.

\textbf{10.} Se $ x $ é uma solução periódica, $ x - x_p = x_h $ é um
solução periódica de $(H)$ que tende a $0$ no $+\infty $; um
função periódica com limite $0$ é idêntico $0$ (seus valores em
um período se repete para sempre, então cada valor é um limite de um
subsequência tendendo a $0$). Por isso $ x = x_p $.

\textbf{11.} Aqui $\lambda = 1$, $\omega_0^2 = 2$, $\Omega = 1$,
$ A = 1$: $ z = \frac1{2 - 1 + 2\iu} = \frac{1 - 2\iu}5$, então
\[
x_p = \Re\Bigl(\frac{(1 - 2\iu)(\cos t + \iu\sin t)}5\Bigr)
= \frac{\cos t + 2\sin t}5 .
\]
Homogêneo: raízes de $ r^2 + 2r + 2$ são $-1 \pm \iu $: $ x_h =
\eu^{-t}(C\cos t + S\sin t)$. Condições: $ x(0) = \frac15 + C = 0$
dá $ C = -\frac15$; diferenciando, $ x'(0) = \frac25 - C + S =
0$ dá $ S = C - \frac25 = -\frac35$. Por isso
\[
x(t) = \underbrace{\frac{\cos t + 2\sin t}5}_{\text{steady}}
- \underbrace{\eu^{-t}\,\frac{\cos t + 3\sin
t}5}_{\text{transient}} ,
\]
o transitório morrendo como $\eu^{-t}$.

\textbf{12.} Expandindo, $ g(u) = u^2 - 2(\omega_0^2 - 2\lambda^2)u
+ \omega_0^4 = (u - u_r)^2 + g(u_r)$ com $ u_r = \omega_0^2 -
2\lambda^2$ e
\[
g(u_r) = \omega_0^4 - u_r^2 = (\omega_0^2 - u_r)(\omega_0^2 +
u_r) = 2\lambda^2\bigl(2\omega_0^2 - 2\lambda^2\bigr) =
4\lambda^2(\omega_0^2 - \lambda^2) .
\]
Se $2\lambda^2 < \omega_0^2$, então $ u_r > 0$ é admissível
frequência quadrada: $ g $ tem um mínimo estrito lá, então $ R = A/\sqrt
g $ tem um máximo estrito em $\Omega_r = \sqrt{u_r} =
\sqrt{\omega_0^2 - 2\lambda^2}$, com $ R_{\max} = A/\sqrt{g(u_r)}
= \frac{A}{2\lambda\sqrt{\omega_0^2 - \lambda^2}}$.

\textbf{13.} $ R(0) = A/\omega_0^2$, então
\[
\frac{R_{\max}}{R(0)}
= \frac{\omega_0^2}{2\lambda\sqrt{\omega_0^2 - \lambda^2}}
= \frac{\omega_0}{2\lambda}\cdot
\frac{\omega_0}{\sqrt{\omega_0^2 - \lambda^2}}
= Q\,\Bigl(1 - \frac{\lambda^2}{\omega_0^2}\Bigr)^{-1/2}
\geq Q .
\]
Para amortecimento fraco o fator de correção é próximo de $1$: \omterm{ex:b1:diffeq:oscillation}{ressonância}
multiplica o deslocamento estático por essencialmente $ Q $.

\textbf{14.} $ V(\Omega)^2 = \frac{A^2 u}{g(u)}$ com $ u =
\Omega^2$. Sua derivada tem o sinal de $ g(u) - u\,g'(u) =
(\omega_0^2 - u)^2 + 4\lambda^2 u - u\bigl(2(u - \omega_0^2) +
4\lambda^2\bigr) = (\omega_0^2 - u)^2 + 2u(\omega_0^2 - u) =
(\omega_0^2 - u)(\omega_0^2 + u)$, positivo para $ u < \omega_0^2$
e negativo além: máximo estrito exatamente em $\Omega =
\omega_0$, para cada amortecimento. Lá $\tan\varphi $ explode com
$\varphi \in \intoo0\pi $: $\varphi = \frac\pi2$, e $ x_p'(t) =
-R\Omega\sin(\Omega t - \frac\pi2) = R\Omega\cos(\Omega t)$ é
exatamente em fase com a força: transferência de potência ideal.

\textbf{15.} $ R = R_{\max}/\sqrt2 \iff g(u) = 2g(u_r) \iff (u -
u_r)^2 = g(u_r) = 4\lambda^2(\omega_0^2 - \lambda^2)$, dando
\[
u_\pm = u_r \pm 2\lambda\sqrt{\omega_0^2 - \lambda^2},
\qquad
\Omega_+^2 - \Omega_-^2 = 4\lambda\sqrt{\omega_0^2 - \lambda^2} .
\]
Então $\Omega_+ - \Omega_- = \frac{\Omega_+^2 -
\Omega_-^2}{\Omega_+ + \Omega_-}$, e para $\lambda \ll \omega_0$
ambos $\Omega_\pm \approx \omega_0$: $\Omega_+ - \Omega_- \approx
\frac{4\lambda\omega_0}{2\omega_0} = 2\lambda $, então
$\frac{\omega_0}{\Omega_+ - \Omega_-} \approx
\frac{\omega_0}{2\lambda} = Q $. Medindo um \omterm{ex:b1:diffeq:oscillation}{ressonância} largura do pico
mede a sua \omterm{pb:b1:diffeq:1}{fator de qualidade}.

\textbf{16.} $ Q = 10$; $\Omega_r = \sqrt{1 - 2(0.05)^2} =
\sqrt{0.995} = 0.9975$; $ R_{\max} = \frac1{2 \times 0.05
\sqrt{1 - 0.0025}} = \frac1{0.1 \times 0.99875} = 10.01$; estático
resposta $ R(0) = 1$; largura de banda $\approx 2\lambda = 0.1$. Um alto
ponta fina de altura $\approx Q $ sobre um planalto de altura $1$.

\textbf{17.} Se $2\lambda^2 \geq \omega_0^2$ então $ u_r \leq 0$ e
$ g'(u) = 2(u - u_r) > 0$ para todos $ u > 0$: $ g $ aumenta estritamente
ligado $\intoo0{+\infty}$, então $ R = A/\sqrt g $ diminui estritamente de
$ R(0) = A/\omega_0^2$: a resposta é maior na frequência zero
e não há pico.

\textbf{18.} $\gamma = \iu\Omega $ não é uma raiz de $ r^2 +
\omega_0^2$ (como $\Omega \neq \omega_0$), então
\cref{met:b1:diffeq:particular} com $ m = 0$ dá $ x_p =
\frac{A\cos\Omega t}{\omega_0^2 - \Omega^2}$ (substituto e
verifique: $-\Omega^2 + \omega_0^2$ vezes o cosseno). Geral
solução:
\[
x(t) = \frac{A\cos\Omega t}{\omega_0^2 - \Omega^2}
+ \lambda_1\cos\omega_0 t + \mu_1\sin\omega_0 t .
\]

\textbf{19.} $ x(0) = 0$ forças $\lambda_1 = -\frac A{\omega_0^2 -
\Omega^2}$, e $ x'(0) = 0$ forças $\mu_1 = 0$:
\[
x(t) = \frac{A}{\omega_0^2 - \Omega^2}\,
\bigl(\cos\Omega t - \cos\omega_0 t\bigr)
= \frac{2A}{\omega_0^2 - \Omega^2}\,
\sin\Bigl(\frac{(\omega_0 - \Omega)t}2\Bigr)
\sin\Bigl(\frac{(\omega_0 + \Omega)t}2\Bigr) ,
\]
pela fórmula do produto $\cos a - \cos b =
2\sin\frac{b - a}2\sin\frac{b + a}2$ aplicado com $ a = \Omega t $,
$ b = \omega_0 t $.

\textbf{20.} Para $\Omega $ aproximar $\omega_0$, o segundo seno
oscila na frequência rápida $\frac{\omega_0 + \Omega}2
\approx \omega_0$, enquanto o primeiro é um envelope lento de frequência
$\frac{\abs{\omega_0 - \Omega}}2$: a amplitude do rápido
oscilação aumenta e diminui com período de envelope
$\frac{2\pi}{\abs{\omega_0 - \Omega}}$ (duas \omterm{pb:b1:diffeq:1}{batidas} por envelope
período), atingindo máximos $\frac{2A}{\abs{\omega_0^2 - \Omega^2}}$.
Como $\Omega \to \omega_0$, as \omterm{pb:b1:diffeq:1}{batidas} tornam-se mais lentas (ponto final
$\to \infty $) e mais alto (amplitude $\to \infty $).

\textbf{21.} Consertar $ t $. Como $\Omega \to \omega_0$:
\[
\frac{2A}{\omega_0^2 - \Omega^2}\sin\Bigl(\frac{(\omega_0 -
\Omega)t}2\Bigr)
= \frac{2A}{\omega_0 + \Omega}\cdot
\frac{\sin\bigl(\frac{(\omega_0 - \Omega)t}2\bigr)}
{\omega_0 - \Omega}
\longrightarrow \frac{2A}{2\omega_0}\cdot\frac t2
= \frac{At}{2\omega_0} ,
\]
enquanto $\sin\bigl(\frac{(\omega_0 + \Omega)t}2\bigr) \to
\sin(\omega_0 t)$: o limite é $ x_\infty(t) =
\frac{At\sin\omega_0 t}{2\omega_0}$. Verificação direta: com $ C =
\frac A{2\omega_0}$, $ x_\infty = Ct\sin\omega_0 t $ tem
$ x_\infty'' = 2C\omega_0\cos\omega_0 t - C\omega_0^2 t
\sin\omega_0 t $, então $ x_\infty'' + \omega_0^2 x_\infty =
2C\omega_0\cos\omega_0 t = A\cos\omega_0 t $, com $ x_\infty(0) =
0$ e $ x_\infty'(0) = C\sin 0 + C\omega_0 \cdot 0 \cdot \cos 0 =
0$ --- para $\omega_0 = A = 1$ isso é
exatamente \cref{ex:b1:diffeq:oscillation}. \omterm{ex:b1:diffeq:oscillation}{Ressonância} é o
degeneração das \omterm{pb:b1:diffeq:1}{batidas}: a primeira expansão do envelope, esticada
para comprimento infinito.

\textbf{22.} Sem amortecimento, a amplitude ressonante cresce linearmente
e sem limite; com amortecimento $\lambda > 0$ o crescimento satura
em $ R_{\max} \approx Q\,\frac A{\omega_0^2}$. O mecanismo:
dissipação remove energia a uma taxa que cresce com a amplitude
(questão 23), então o acúmulo para exatamente quando o amortecimento
queima energia tão rápido quanto o forçamento a fornece.

\textbf{23.} Com $ x_p' = -R\Omega\sin(\Omega t - \varphi)$, acabou
um período as médias $\langle\cos^2\rangle =
\langle\sin^2\rangle = \frac12$ e $\langle\sin\cos\rangle = 0$
dar:
\[
\langle P_{\mathrm{diss}}\rangle
= 2\lambda\,R^2\Omega^2\,\langle\sin^2(\Omega t -
\varphi)\rangle = \lambda R^2\Omega^2 ;
\]
e expandindo $\sin(\Omega t - \varphi) = \sin\Omega t\cos\varphi
- \cos\Omega t\sin\varphi $:
\[
\langle P_{\mathrm{in}}\rangle
= -AR\Omega\,\bigl\langle\cos\Omega t\,\sin(\Omega t -
\varphi)\bigr\rangle
= AR\Omega\,\frac{\sin\varphi}2 .
\]
Desde $\sin\varphi = \frac{2\lambda\Omega}{\sqrt D} =
\frac{2\lambda\Omega R}A $, isso é $\frac{AR\Omega}2 \cdot
\frac{2\lambda\Omega R}A = \lambda R^2\Omega^2$: injetado e
os poderes dissipados se equilibram exatamente --- a propriedade definidora de um
regime constante.

\textbf{24.} (i) O teorema da estrutura divide cada solução em
estado estacionário mais transitório (questões 9 a 11) e reduzido
singularidade para o problema homogêneo. (ii) O método complexo
transformou a busca por uma solução específica em uma \omterm{thm:b1:logic:partition}{divisão} de
números complexos (questão 6), com amplitude e fase lidas em um
\omterm{def:b1:complex:field}{módulo} e um argumento. (iii) O \omterm{ex:b1:diffeq:oscillation}{ressonância} curva é pura
estudo de função --- uma quadrática em $ u = \Omega^2$, seu mínimo,
seu nível conjuntos --- no estilo de \cref{ch:b1:functions}
(questões 12 a 17).

\textbf{25.} Livre e amortecido: pseudo-oscilações decadentes,
vida inteira $\frac1\lambda $, sobre $\frac Q\pi $ balanços. Forçado e
amortecido: os transientes morrem e um estado estacionário sinusoidal único
sobrevive na frequência de forçamento, com pico de amplitude próximo
$\omega_0$ (altura $\approx Q \times $ estático, largura $\approx
\frac{\omega_0}Q $) e varredura de fase de $0$ a $\pi $ através
$\frac\pi2$ no $\omega_0$. Livre e não amortecido: perpétuo
oscilação. Forçado e não amortecido: \omterm{pb:b1:diffeq:1}{batidas}, degenerando em
crescendo linearmente \omterm{ex:b1:diffeq:oscillation}{ressonância} na sintonia exata. Um adimensional
número, $ Q = \frac{\omega_0}{2\lambda}$, sintoniza tudo ---
altura de pico, largura de banda e vida útil transitória são três leituras
do mesmo mostrador. A sequência: reescrita $ x'' + 2\lambda x' +
\omega_0^2x$ como um sistema de primeira ordem abre os métodos matriciais de
\cref{ch:b1:matrices} e o volume do Ano 2, e decompondo um
forçamento periódico arbitrário em sinusóides (série de Fourier, Ano 3
volume) torna a análise de frequência única deste problema o
bloco de construção universal: resolva para cada frequência, sobreponha.
\end{solution}
